\chapter{Exam Blueprint}
\section{Domains and Weights}
\begin{itemize}[leftmargin=*]
  \item Composition — 23\%
  \item Content Aggregation — 10\%
  \item Customizing USD — 6\%
  \item Data Exchange — 15\%
  \item Data Modeling — 13\%
  \item Debugging and Troubleshooting — 11\%
  \item Pipeline Development — 14\%
  \item Visualization — 8\%
\end{itemize}

\section{How to study (fast feedback loop)}
\begin{itemize}[leftmargin=*]
  \item Read cheat sheet $\rightarrow$ do drills $\rightarrow$ explain to yourself $\rightarrow$ re-run drills.
  \item Prefer tiny USD layers you can reason about (3–20 lines).
  \item Always ask: “Which layer wins and why?” (LIVRPS + strength).
\end{itemize}

\section{Field notes from test-takers}
Synthesized from public write-ups by certified developers and from this
author's own attempts --- the things people consistently wish they had known
before booking the exam.

\begin{itemize}[leftmargin=*]
  \item \textbf{Experience does not equal readiness.} Developers with years
        of production USD and Omniverse work report failing a first attempt
        after $\sim$80 hours of study. The exam tests the guide's objectives
        \emph{specifically} --- including domains your daily work never
        touches.
  \item \textbf{Budget 80--120 focused hours} over 8--10 weeks. People who
        passed recommend at least 100; under-investment is the most-cited
        cause of first-attempt failures.
  \item \textbf{Work the guide point by point.} Coverage beats extra depth in
        your comfort zone. The margin tags and the Coverage Map appendix in
        this book exist for exactly that: no objective left unvisited.
  \item \textbf{Composition dominates --- mentally.} Expect LIVRPS scenarios
        you must trace \emph{without} usdview. Practice predicting the
        composed result first and verifying second; then practice with the
        verification switched off.
  \item \textbf{It is a Python exam.} API questions use the Python surface;
        read and write \texttt{pxr} code fluently.
  \item \textbf{Understand the wrong answers.} For every practice question,
        articulate why each distractor fails --- this book's answer key is
        written that way on purpose. Then build your own variant of the
        scenario; the drills exist for exactly this.
\end{itemize}

\noindent\textbf{Logistics.} The exam is remotely proctored (Certiverse),
costs about \$200, and the certification is valid for two years. A failed
attempt can be retaken after 14 days, up to five attempts per year --- plan
your calendar so a retake is an option, not a crisis. For hands-on practice,
\texttt{usdview} alone is sufficient tooling; no special hardware is needed.
